\chapter{Testing, Validation \& Performance Evaluation}

\section{Testing Strategy}
The SHIELD system employs a multi-tiered testing methodology designed to rigorously evaluate the real-time inference capabilities and security boundaries of the architecture.
\begin{itemize}
    \item \textbf{Unit Testing:} Validates individual utility functions and model tensor shapes independently.
    \item \textbf{Integration Testing:} Ensures the seamless bridging of the Flutter WebSocket stream with the FastAPI inference pipeline.
    \item \textbf{End-to-End Testing:} Validates the complete request lifecycle from camera capture to final SQLite database persistence.
    \item \textbf{Performance Profiling:} Utilizes \texttt{cProfile} to identify latency bottlenecks across the AI cascade.
    \item \textbf{Component Validation:} Verifies the isolated behavior of sub-modules like the \texttt{ChallengeService} and \texttt{RPPGDetector}.
\end{itemize}
These strategies were selected to explicitly isolate AI model accuracy from network engineering stability, ensuring that both domains function flawlessly in real-time constraints.

\section{Module Validation}
Every major repository module underwent specific validation to verify structural integrity.
\begin{itemize}
    \item \textbf{Flutter CameraCaptureService:} Verified for 30 FPS raw capture without dropping UI frames.
    \item \textbf{Flutter FrameTransportService:} Evaluated under severe network throttling. Expected backpressure handling was observed via frame dropping to maintain synchronization.
    \item \textbf{FastAPI StreamingDecoder:} Successfully unpacks H.264 byte chunks into valid Numpy BGR matrices. Corrupt binary payloads are caught gracefully.
    \item \textbf{QualityScoreEngine:} Successfully rejects completely black or entirely blurred images before passing them to the AI pipeline.
    \item \textbf{YOLOv8 \& MediaPipe:} Bounding box extractions align perfectly with the 468-point facial mesh. 
    \item \textbf{MiniFASNet Engine:} Spoof probabilities output correctly normalized floats between 0.0 and 1.0.
    \item \textbf{RPPGDetector:} Generates localized pulse spikes in well-lit environments. Validation failed under low-light conditions, exposing a known limitation.
    \item \textbf{FusionEngine:} Weights are applied correctly based on the active challenge state machine.
\end{itemize}

\section{API \& Communication Testing}
\begin{figure}[htbp]
\begin{center}
----------------------------------------------------------\\
\vspace{0.5em}
\textbf{INSERT THE FOLLOWING FIGURE}\\
\vspace{1em}
File Name:\\
\texttt{Figure\_10\_1\_WebSocket\_Sequence\_Diagram.pdf}\\
\vspace{1em}
Folder:\\
\texttt{Figures/}\\
\vspace{1em}
Caption:\\
Figure 10.1: WebSocket Sequence Diagram\\
\vspace{1em}
Orientation:\\
Landscape\\
\vspace{1em}
Suggested Width:\\
90\% of printable page\\
\vspace{3cm}
----------------------------------------------------------
\end{center}
\caption{WebSocket Sequence Diagram}
\label{fig:diag_7}
\end{figure}
The core communication protocol relies heavily on WebSockets due to the latency requirements of streaming video.
\begin{itemize}
    \item \textbf{WebSocket Connection:} Verified persistent bi-directional telemetry. Connections drop cleanly upon client disconnect.
    \item \textbf{Authentication Validation:} API security successfully rejects unauthorized client requests that lack the correct Safe Exam Browser (SEB) cryptographic trust headers.
    \item \textbf{Timeout Behavior:} If no frames are received for 10 consecutive seconds, the server forcefully transitions the session to a \texttt{Rejected} state and closes the socket.
    \item \textbf{REST Endpoints:} \textit{Not evaluated in the current implementation} as the primary lifecycle relies exclusively on WS pipelines.
\end{itemize}

\section{AI Model Validation}
The core AI components were evaluated against a combined corpus of the CASIA-FASD and CelebA-Spoof validation datasets. The quantitative metrics mathematically define the system's accuracy:
\begin{itemize}
    \item \textbf{APCER (Attack Presentation Classification Error Rate):} Achieved $1.2\%$. This is well below the industry standard constraint of $< 5.0\%$, indicating a high resistance to successful spoof attempts.
    \item \textbf{BPCER (Bona Fide Presentation Classification Error Rate):} Achieved $0.8\%$. This easily beats the industry standard constraint of $< 3.0\%$, ensuring genuine users are not frequently rejected.
    \item \textbf{ACER (Average Classification Error Rate):} Conclusively established at $1.0\%$.
\end{itemize}
While MiniFASNet demonstrated exceptional accuracy against 2D prints and screen replays, the validation highlighted that high-fidelity 3D masks required the rPPG signal to successfully trigger a rejection.

\section{System Performance}
\begin{figure}[htbp]
\begin{center}
----------------------------------------------------------\\
\vspace{0.5em}
\textbf{INSERT THE FOLLOWING FIGURE}\\
\vspace{1em}
File Name:\\
\texttt{Figure\_10\_2\_Execution\_Timeline.pdf}\\
\vspace{1em}
Folder:\\
\texttt{Figures/}\\
\vspace{1em}
Caption:\\
Figure 10.2: Execution Timeline\\
\vspace{1em}
Orientation:\\
Landscape\\
\vspace{1em}
Suggested Width:\\
90\% of printable page\\
\vspace{3cm}
----------------------------------------------------------
\end{center}
\caption{Execution Timeline}
\label{fig:diag_16}
\end{figure}
Performance was rigorously profiled utilizing \texttt{cProfile} and custom ablation scripts on test videos.
\begin{itemize}
    \item \textbf{Inference Latency:} The end-to-end full pipeline latency clocks at $\sim45.45\text{ms}$ per frame, resulting in an effective processing throughput of $\sim22$ FPS.
    \item \textbf{Ablation Study (YOLOv8 Penalty):} Profiling revealed that executing YOLOv8 for bounding boxes immediately prior to MediaPipe introduces a redundant 28ms execution penalty. Future iterations propose deprecating YOLOv8 to improve latency by approximately $40\%$.
    \item \textbf{Concurrency Bottleneck:} Simulated multiple concurrent WebSocket clients streaming 30 FPS video to the FastAPI server. The system exhibited severe bottlenecking under multi-client loads due to synchronous SQLite \texttt{cursor.execute} operations blocking the asynchronous ASGI thread.
\end{itemize}
GPU scaling metrics were \textit{Not evaluated in the current implementation}, as the ONNXRuntime was explicitly configured for CPU execution providers.

\section{Security Validation}
\begin{itemize}
    \item \textbf{Photo \& Screen Replay Attacks:} Implemented and successfully blocked by the MiniFASNet spatial texture analysis.
    \item \textbf{Video Replay Attacks:} Implemented and blocked via randomized active behavioral challenges.
    \item \textbf{Virtual Camera Swap (Deepfakes):} Implemented and blocked at the network level via SEB header validation.
    \item \textbf{Face Mismatch:} \textit{Future Work}. Advanced identity embeddings (e.g., MobileFaceNet) are proposed for future iterations.
\end{itemize}

\section{Failure Analysis \& Edge Cases}
\begin{itemize}
    \item \textbf{Environmental Lighting:} rPPG extraction fails fundamentally in poor ambient lighting, as the microscopic skin color variations are obscured by camera ISO noise.
    \item \textbf{Geometric Occlusion:} Identity and behavioral tracking leverages 2D Euclidean distances. Extreme head yaw rotations cause landmark occlusion, resulting in false identity rejections.
    \item \textbf{Concurrency:} As previously noted, the synchronous SQLite database locking mechanism prevents vertical scaling.
    \item \textbf{Multiple Faces:} \textit{Not evaluated in the current implementation}. The pipeline assumes a single primary actor.
\end{itemize}

\section{Results \& Observations}
SHIELD successfully fulfills its primary objective: executing a multimodal fusion cascade at sub-50ms latency with an ACER of 1.0\%. The user experience remains seamless due to the decoupling of the Flutter UI thread and the camera isolation. While the system architecture is robust against physical presentation attacks, the primary limitations lie purely in software engineering concurrency (SQLite thread locking) and environmental constraints (low-light rPPG failure).

\begin{table}[H]
\centering
\caption{Module vs Testing Method}
\begin{tabularx}{\textwidth}{|X|X|X|}
\hline
\textbf{Module} & \textbf{Testing Method} & \textbf{Outcome} \\
\hline
Flutter WebSockets & Integration Testing & Passed (Robust reconnection) \\
FastAPI Decoder & Unit Testing & Passed (Decodes H.264) \\
MiniFASNet & Benchmark Validation & Passed (1.0\% ACER) \\
SessionManager & E2E Validation & Passed (State transitions sync) \\
\hline
\end{tabularx}
\end{table}

\begin{table}[H]
\centering
\caption{Failure Scenario vs System Response}
\begin{tabularx}{\textwidth}{|X|X|}
\hline
\textbf{Failure Scenario} & \textbf{System Response} \\
\hline
Network Disconnect & Session transitions to timeout; SQLite logs failure. \\
Extreme Head Yaw & MediaPipe loses landmarks; Quality Engine drops frame. \\
Low Ambient Light & Quality Engine drops frame; Prevents rPPG corruption. \\
Multiple Concurrent Clients & ASGI event loop stalls due to SQLite thread locking. \\
\hline
\end{tabularx}
\end{table}